\documentclass[10pt,a4paper]{amsart}
\usepackage[margin=19mm]{geometry}
\usepackage{amsmath,amssymb,mathtools}
\usepackage[scr=rsfs,cal=euler]{mathalfa}
\usepackage{xfrac}
\usepackage[unicode,hidelinks,bookmarks=false]{hyperref}
\newcommand{\ie}{i.e.\ }
\newcommand{\cf}{cf.\ }
\newcommand{\resp}{respectively\ }
\newcommand{\Subsection}{\subsection}
\providecommand{\nobreakdash}{\nobreak\hbox{-}}
\setlength{\emergencystretch}{2em}
\multlinegap0pt
\usepackage{bm}
\usepackage{bbm}
\usepackage{dsfont}
\usepackage{xspace}
\usepackage{cancel}
\usepackage{mathtools}
\usepackage{cleveref}
\usepackage{amsthm}
\makeatletter
\@ifundefined{theorem}{\newtheorem{theorem}{Theorem}}{}
\@ifundefined{proposition}{\newtheorem{proposition}[theorem]{Proposition}}{}
\makeatother
\newcommand{\bkt}[1]{\left[#1\right]}
\newcommand{\eqdef}{\vcentcolon=}
\newcommand{\mbbE}{\mathbb{E}}
\newcommand{\mbbP}{\mathbb{P}}
\newcommand{\mbbR}{\mathbb{R}}
\newcommand{\mbb}[1]{\mathbbm{#1}}
\newcommand{\mcT}{\mathcal{T}}
\newcommand{\mcX}{\mathcal{X}}
\newcommand{\norm}[1]{\left\lVert#1\,\right\rVert}
\title[Asymptotically optimal reinforcement learning in Block Marko]{Asymptotically optimal reinforcement learning in Block Markov Decision Processes}
\author{Thomas van Vuren, Fiona Sloothaak, Maarten G. Wolf, Jaron Sanders}
\date{Source: arXiv:2510.13748v1 (CC BY 4.0)}
\begin{document}
\maketitle

\noindent\emph{Source and licence.} Asymptotically optimal reinforcement learning in Block Markov Decision Processes. Version arXiv:2510.13748v1 (CC BY 4.0), distributed under the Creative Commons Attribution 4.0 International licence, as recorded in the arXiv licence metadata for this exact version and shown on the abstract page https://arxiv.org/abs/2510.13748v1. Full rights notice: licenses/arxiv-2510.13748-CC-BY-4.0.txt.\par\smallskip

\noindent\textit{Editorial scope.} Counted unit: the Proposition whose source label is \texttt{prop:concentration\_inequalities} in the article (source0.tex lines 6672--6724, source label \texttt{prop:concentration\_inequalities}), delivered together with the article's own complete original proof of it (source0.tex lines 6772--6997, one \texttt{proof} environment of 226 lines). No mathematics is written, repaired or re-derived: statement and proof are the article's own text line for line. The proof is detached from the statement in the source: the article states the result at lines 6672--6724 and prints its proof at lines 6772--6997, and the proof environment's own header names the statement, so the association is the article's own. Required context retained uncounted: eqn:countsdef\_overload (lines 1534--1538); eqn:pfqfdef (lines 6658--6669). Every definition, display and label of those lines is retained unchanged. Omitted from this package and not redistributed: 1--1533, 1539--6657, 6670--6671, 6725--6771, 6998--8069. Nothing inside a retained excerpt is omitted or altered: every retained body is delivered exactly as the article's lines are written, so no omission receipt of any kind accompanies this package. Citations: the retained excerpts carry no citation command at all, so no bibliography entry of the article is transcribed and no reference is redistributed. Reference adaptation: none was needed and none was made; every reference inside the delivered unit resolves inside it, so no unresolved reference or citation is produced. Printed slips kept literal: no printed word, symbol, hypothesis or number of the counted statement or of its proof was altered. Layout and class: the article's own class is \texttt{article} with packages including geometry, inputenc, amsmath, amssymb, amsthm, authblk, bbm, comment, enumitem, mathtools, subcaption, biblatex, tikz, pgfplots; the wrapper is the lane's pinned \texttt{amsart} editorial wrapper at 10pt on A4, which restates the theorem-like environments the retained text needs and adds the article's own macro definitions that the retained text needs (\texttt{\textbackslash bkt}, \texttt{\textbackslash eqdef}, \texttt{\textbackslash mbb}, \texttt{\textbackslash mbbE}, \texttt{\textbackslash mbbP}, \texttt{\textbackslash mbbR}, \texttt{\textbackslash mcT}, \texttt{\textbackslash mcX}, \texttt{\textbackslash norm}), with extra wrapper packages (bm, bbm, dsfont, xspace, cancel, mathtools, cleveref). The delivered numbering is the unit's own. Assets: no retained span contains a figure environment, a graphics-inclusion command, a PDF-page inclusion, a table or a TikZ picture; no image file is copied into this package, the package declares no asset dependency and no image file is redistributed with it. Licence: version arXiv:2510.13748v1 of this article is distributed under the Creative Commons Attribution 4.0 International licence, as recorded in the arXiv licence metadata for this exact version and shown on the abstract page https://arxiv.org/abs/2510.13748v1. A full rights notice is delivered at licenses/arxiv-2510.13748-CC-BY-4.0.txt. Characters: every retained excerpt is pure ASCII, so the delivered engine, XeLaTeX with its default Latin Modern text font, reports no missing character.
\begin{equation}
	\label{eqn:countsdef_overload}
	\hat{N}_{k}(\mcX_1,a)\eqdef \hat{N}_{k}(\mcX_1,a,[n])\quad \mathrm{and} \quad
	\hat{N}_{k}^{\mathrm{to}}(\mcX_2)\eqdef \sum_{a}\hat{N}_{k}([n],a,\mcX_2).
\end{equation}

\begin{equation}
	\label{eqn:pfqfdef}
	\hat{p}_k^f(s'\mid s,a)
	\eqdef
	\frac{\hat{N}_k(f^{-1}(s),a,f^{-1}(s'))}{1\vee \hat{N}_k(f^{-1}(s),a)}
	,
	\!\quad \mathrm{and},\! \quad
	\hat{q}_k^f(x\mid s)
	\eqdef
	\frac{\mbb{1}\{x\in f^{-1}(s)\}\hat{N}^{\mathrm{to}}_k(x)}{1\vee \hat{N}^{\mathrm{to}}_k(f^{-1}(s))}
	.
\end{equation}

\begin{proposition}
	\label{prop:concentration_inequalities}
	Let $U\in\mbbR^n_{\geq0}$ and $u\in\mbbR^S_{\geq0}$ be two vectors with nonnegative entries, and let $\delta\in(0,1/2]$.
	Then the following inequalities hold for $s,s'\in[S]$, $a\in[A]$, $y\in[n]$:
	\begin{align}
		 &
		\mbbP\biggl[
		|\langle \hat{q}_k^f(\,\cdot \mid s)-q(\,\cdot \mid s), U\rangle|
		>
		\sqrt{\frac{\lVert U\rVert_{\infty}^2 \ln(2 T_k/\delta)}{1\vee \hat{N}^{\mathrm{to}}_k(f^{-1}(s))}}
		\biggr]
		\leq
		\delta
		\label{eqn:hoeffding_q}
		\\
		 &
		\mbbP\biggl[
			|\langle \hat{p}_k^f(\,\cdot \mid s,a)-p(\,\cdot \mid s,a), u\rangle|
			>
			\sqrt{\frac{\lVert u\rVert_{\infty}^2\ln(2 T_k/\delta)}{1\vee \hat{N}_k(f^{-1}(s),a)}}
			\biggr]
		\leq
		\delta
		,
		\label{eqn:hoeffding_p}
		\\
		 &
		\mbbP\biggl[
		|\hat{q}_k^f(y\mid s)-q(y\mid s)|
		>
		\sqrt{\frac{q(y\mid s) \ln(2T_k/\delta)}{1\vee \hat{N}^{\mathrm{to}}_k(f^{-1}(s))}}
		+
		\frac{2\ln(2T_k/\delta)}{3(1\vee\hat{N}^{\mathrm{to}}_k(f^{-1}(s)))}
		\biggr]
		\leq
		\delta
		,
		\label{eqn:bernstein_q}
		\\
		 &
		\mbbP\biggl[
			| \hat{p}_k^f(s'\mid s,a)-p(s'\mid s,a)|
			>
			\sqrt{\frac{p(s'\mid s,a)\ln(2T_k/\delta)}{1\vee \hat{N}_k(f^{-1}(s),a)}}
			+
			\frac{2\ln(2T_k/\delta)}{3(1\vee\hat{N}_k(f^{-1}(s),a))}
			\biggr]
		\leq
		\delta
		.
		\label{eqn:bernstein_p}
	\end{align}
\end{proposition}

\begin{proof}[Proof of \Cref{prop:concentration_inequalities}]
	We first prove \eqref{eqn:hoeffding_q}.
	Fix $s$ and $U\in\mbbR^n_{\geq0}$, and let $\delta\in(0,1/2]$.

	Let $\mcT'_k(s) \eqdef \{(l,h)\in [k-1]\times [H]: x_{l,h+1}\in f^{-1}(s)\}$ denote the (random) set of times for which $f(x_{l,h+1})=s$.
	Note that $(i)$ $|\mcT'_k(s)|=\hat{N}_k^{\mathrm{to}}(f^{-1}(s))$.
	Conditional on $\{\hat{N}_k^{\mathrm{to}}(f^{-1}(s))=N\}$ for $N>0$ it follows from \eqref{eqn:pfqfdef} and \eqref{eqn:countsdef_overload} that
	\begin{align}
		 &
		N\langle \hat{q}^f_k(\,\cdot \mid s)-q(\,\cdot \mid s), U\rangle
		=
		\sum_{y\in [n]}
		\biggl(
		\mathbbm{1}\{y\in f^{-1}(s)\}\sum_{l\in[k-1]}\sum_{h=1}^H\mathbbm{1}\{x_{l,h+1}=y\}-Nq(y\mid s)
		\biggr)
		U(y)
		\nonumber
		\\
		 &
		=
		\sum_{y\in [n]}
		\biggl(\sum_{l\in[k-1]}\sum_{h=1}^H
		\mathbbm{1}\{x_{l,h+1}\in f^{-1}(s)\}\mathbbm{1}\{x_{l,h+1}=y\}-Nq(y\mid s)
		\biggr)
		U(y)
		\nonumber
		\\
		 &
		\overset{(i)}{=}
		\sum_{l\in[k-1]}
		\sum_{h=1}^H
		\mathbbm{1}\{x_{l,h+1} \in f^{-1}(s)\}\sum_{y\in [n]}(
		\mathbbm{1}\{x_{l,h+1}=y\}-q(y\mid s)
		)
		U(y)
		\nonumber
		\\
		 &
		=
		\sum_{(l,h)\in \mcT'_k(s)}
		\bigl(U(x_{l,h+1})-\sum_{y\in f^{-1}(s)}q(y\mid s) U(y)\bigr),
		\label{eqn:q_esterror_distributed_as}
	\end{align}
	where in passing to the final line we used that $q(y\mid s)=0$ unless $y\in f^{-1}(s)$.

	By definition of $\mcT'_k(s)$ and the Markov property, each $x_{l,h+1}$ for $(l,h)\in \mcT'_k(s)$ is independent and identically distributed according to $q(\,\cdot \mid s)$.
	It then follows from \eqref{eqn:q_esterror_distributed_as} that conditional on $\{\hat{N}^{\mathrm{to}}_k(f^{-1}(s))=N\}$ for $N>0$, $N\langle \hat{q}^f_k(\,\cdot \mid s)-q(\,\cdot \mid s), U\rangle$ is the sum of $N$ i.i.d. random variables $X_t$ for $t=1,\ldots,N$ with distribution $\mbbP[X_t = U(x)-\sum_{y\in f^{-1}(s)}q(y\mid s) U(y)]=q(x\mid s)$ for $x\in[n]$, from which one also checks that $\mbbE[X_t] = 0$.
	Because $U$ has nonnegative entries, the $X_t$ are almost surely bounded in $[0,\lVert U\rVert_{\infty}]$.

	It then follows from Hoeffding's inequality that for any $N>0$ and $\epsilon>0$
	\begin{equation}
		\label{eqn:basic_hoeffding}
		\mbbP\bkt{
		|\langle \hat{q}^f_k(\,\cdot \mid s)-q(\,\cdot \mid s), U\rangle| > \epsilon
		\mid
		\hat{N}^{\mathrm{to}}_k(f^{-1}(s))=N
		}
		\leq
		2\exp\biggl(
		-
		\frac{2N\epsilon^2}{\lVert U\rVert_{\infty}^2}
		\biggr)
		.
	\end{equation}
	Moreover, when $\hat{N}^{\mathrm{to}}_k(f^{-1}(s))=0$, the definition of $\hat{q}_k^f$ in \eqref{eqn:pfqfdef} implies that
	\begin{equation}
		\label{eqn:case_N_equals_zero_hoeffding}
		|\langle \hat{q}_k^f(\,\cdot \mid s) - q(\,\cdot \mid s), U\rangle |
		=
		|\langle q(\,\cdot \mid s), U\rangle|
		\leq
		\lVert U\rVert_{\infty}
		\leq
		\sqrt{\lVert U\rVert_{\infty}^2\ln(2T_k/\delta)}
	\end{equation}
	almost surely for $\delta<1/2$, where we have used that $q(y\mid s)\geq 0$ and $\sum_{y}q(y\mid s)=1$.
	We then rewrite the left-hand side of \eqref{eqn:hoeffding_q} by conditioning on all possible values of $\hat{N}^{\mathrm{to}}_k(f^{-1}(s))$, which is bounded almost surely by $T_k$, and using \eqref{eqn:case_N_equals_zero_hoeffding} and $\mbbP[\hat{N}^{\mathrm{to}}_k(f^{-1}(s))=N]\leq 1$:
	\begin{equation}
		\label{eqn:sum_over_all_values_of_N}
		\begin{split}
			 &
			\mbbP\biggl[
			|\langle \hat{q}^f_k(\,\cdot \mid s)-q(\,\cdot \mid s), U\rangle |>\sqrt{\frac{\norm{U}_{\infty}^2\ln(2T_k/\delta)}{1\vee\hat{N}^{\mathrm{to}}_k(f^{-1}(s))}}
			\biggr]
			\\
			 &
			\leq
			\sum_{N=1}^{T_k}
			\mbbP\biggl[
			|\langle \hat{q}^f_k(\,\cdot \mid s)-q(\,\cdot \mid s), U\rangle|>\sqrt{\frac{\norm{U}_{\infty}^2\ln(2T_k/\delta)}{N}}
			\,\bigg|\,
			\hat{N}^{\mathrm{to}}_k(f^{-1}(s)) = N
			\biggr]
			.
		\end{split}
	\end{equation}
	Finally, applying \eqref{eqn:basic_hoeffding} with $\epsilon = \sqrt{\smash[b]{\norm{U}_{\infty}^2\ln(2T_k/\delta)/N}}$, we conclude that
	\begin{equation}
		\mbbP\biggl[
		|\langle \hat{q}^f_k(\,\cdot \mid s)-q(\,\cdot \mid s), U\rangle |>\sqrt{\frac{\norm{U}_{\infty}^2\ln(2T_k/\delta)}{1\vee\hat{N}^{\mathrm{to}}_k(f^{-1}(s))}}
		\biggr]
		\leq
		2\sum_{N=1}^{T_k}
		\exp(
		-
		\ln(2T_k/\delta)
		)
		=
		2\sum_{N=1}^{T_k}\frac{\delta}{2T_k}
		=
		\delta.
	\end{equation}
	This concludes the proof of \eqref{eqn:hoeffding_q}.

	We next prove \eqref{eqn:bernstein_q}.
	Fix $s$ and $\delta \in(0,1/2]$.
	Note that \eqref{eqn:bernstein_q} holds trivially for $y\notin f^{-1}(s)$ since then $\hat{q}_k^f(y\mid s) = q(y\mid s) = 0$ by definition.
	Hence, fix $y\in f^{-1}(s)$.

	Observe that $\hat{q}_k^f(y\mid s)-q(y\mid s)=\langle\hat{q}_k^f(\,\cdot \mid s)-q(\,\cdot \mid s),e_y\rangle$ with $e_y\in\mbbR^n$ the $y$'th unit vector whose entries are given by $e_y(x)=\mathbbm{1}\{x=y\}$ for $x\in[n]$.
	It then follows from the argument leading to \eqref{eqn:basic_hoeffding} that conditional on $\{\hat{N}^{\mathrm{to}}_k(f^{-1}(s))=N\}$ for $N>0$, $N(\hat{q}^f_k(y\mid s)-q(y\mid s))$ is the sum of $N$ i.i.d. random variables $Y_t$ for $t=1,\ldots,N$ taking values in $\{-q(y\mid s),1-q(y\mid s)\}$ and with distribution $\mbbP[Y_t=1-q(y\mid s)]=q(y\mid s)$ and $\mbbP[Y_t = -q(y\mid s)]=1-q(y\mid s)$.
	It follows that $|Y_t|\leq 1$ almost surely, $\mbbE[Y_t]=0$, and that $Y_t$ has variance bounded by $(1-q(y\mid s))q(y\mid s)\leq q(y\mid s)$.

	We then conclude using Bernstein's inequality that for any $N>0$ and $\epsilon>0$,
	\begin{equation}
		\label{eqn:basic_bernstein}
		\mbbP\bkt{
		| \hat{q}^f_k(y\mid s)-q(y\mid s)| > \epsilon
		\mid
		\hat{N}^{\mathrm{to}}_k(f^{-1}(s))=N
		}
		\leq
		2\exp\biggl(
		-
		\frac{\frac{1}{2}N\epsilon^2}{q(y\mid s)+\frac{1}{3}\epsilon}
		\biggr)
		.
	\end{equation}
	Moreover, let $\epsilon=\epsilon^*(N,L)$ for $N,L>0$ denote the positive solution to $(1/2)N\epsilon^2=L(q(y\mid s)+(1/3)\epsilon)$, and observe that from the elementary inequality $\sqrt{a+b}\leq \sqrt{a}+\sqrt{b}$ for $a,b>0$,
	\begin{equation}
		\label{eqn:optimal_bernstein_epsilon}
		\epsilon^*(N,L)
		=
		\frac{1}{3N}
		\bigl(L + L\sqrt{L+18Nq(y\mid s)}\bigr)
		\leq
		\sqrt{\frac{2q(y\mid s)L}{N}} + \frac{2L}{3N}
		.
	\end{equation}
	Finally, observe also that when $\hat{N}^{\mathrm{to}}_k(f^{-1}(s))=0$, the definition of $\hat{q}_k^f$ in \eqref{eqn:pfqfdef} implies that for $\delta\leq 1/2$ and $T_k\geq 1$,
	\begin{equation}
		\label{eqn:bernstein_case_N_0}
		|\hat{q}_k^f(y\mid s)-q(y\mid s)|
		=
		q(y\mid s)
		\leq
		\sqrt{q(y\mid s)}
		\leq
		\sqrt{q(y\mid s)\ln(2T_k/\delta)}+(2/3)\ln(2T_k/\delta)
		,
	\end{equation}
	where we have used that $q(y\mid s)\in[0,1]$.
	We then rewrite the left-hand side of \eqref{eqn:bernstein_q} in a manner analogous to that in \eqref{eqn:sum_over_all_values_of_N}, but using \eqref{eqn:bernstein_case_N_0} instead of \eqref{eqn:case_N_equals_zero_hoeffding}, and subsequently using \eqref{eqn:optimal_bernstein_epsilon}:
	\begin{equation}
		\begin{split}
			 &
			\mbbP\biggl[
			| \hat{q}^f_k(y\mid s)-q(y\mid s) | >
			\sqrt{\frac{q(y\mid s) \ln(2T_k/\delta)}{1\vee\hat{N}^{\mathrm{to}}_k(f^{-1}(s))}}
			+
			\frac{2\ln(2T_k/\delta)}{3(1\vee\hat{N}^{\mathrm{to}}_k(f^{-1}(s)))}
			\biggr]
			\\
			 &
			\leq
			\sum_{N=1}^{T_k}
			\mbbP\biggl[
			|\hat{q}^f_k(y\mid s)-q(y\mid s)|>
			\sqrt{\frac{q(y\mid s) \ln(2T_k/\delta)}{N}}
			+
			\frac{2\ln(2T_k/\delta)}{3N}
			\,\bigg|\,
			\hat{N}^{\mathrm{to}}_k(f^{-1}(s)) = N
			\biggr]
			\\
			 &
			\leq
			\sum_{N=1}^{T_k}
			\mbbP\bigl[
			|\hat{q}^f_k(y\mid s)-q(y\mid s)|>
			\epsilon^*(N,\ln(2T_k/\delta))
			\mid
			\hat{N}^{\mathrm{to}}_k(f^{-1}(s)) = N
			\bigr]
			.
		\end{split}
	\end{equation}
	It finally follows by applying \eqref{eqn:basic_bernstein} with $\epsilon=\epsilon^*(N,\ln(2T_k/\delta))$ that
	\begin{equation}
		\begin{split}
			\mbbP\biggl[
			| \hat{q}^f_k(y\mid s)-q(y\mid s) |
			>
			 &
			\sqrt{\frac{q(y\mid s) \ln(2T_k/\delta)}{1\vee\hat{N}^{\mathrm{to}}_k(f^{-1}(s))}}
			+
			\frac{2\ln(2T_k/\delta)}{3(1\vee\hat{N}^{\mathrm{to}}_k(f^{-1}(s)))}
			\biggr]
			\\
			 &
			\leq
			2\sum_{N=1}^{T_k}
			\exp(
			-
			\ln(2T_k/\delta)
			)
			=
			2\sum_{N=1}^{T_k}\frac{\delta}{2T_k}
			=
			\delta.
		\end{split}
	\end{equation}
	This concludes the proof of \eqref{eqn:bernstein_q}.

	The analogs \eqref{eqn:hoeffding_p} and \eqref{eqn:bernstein_p} for $\hat{p}_k^f$ of \eqref{eqn:hoeffding_q} and \eqref{eqn:bernstein_q} follow from exactly the same argument, where we replace $\mcT'_k(s)$ by $\mcT'_k(s,a)=\mcT'_k(s,a)\eqdef \{(l,h)\in[k-1]\times [H]:x_{l,h}\in f^{-1}(s),a_{l,h}=a\}$ and condition on $\{\hat{N}_k(f^{-1}(s),a)=N\}$ instead.
\end{proof}

\end{document}
